% Self-contained uncompiled TeX research source.
% No external inputs or bibliography files are required.
% This source has not been TeX-compiled or visually rendered.
\documentclass[11pt]{article}
\usepackage[T1]{fontenc}
\usepackage[utf8]{inputenc}
\usepackage[margin=1in]{geometry}
\usepackage{amsmath,amssymb,amsthm,mathtools,enumitem}
\usepackage[hidelinks]{hyperref}
\newtheorem{theorem}{Theorem}[section]
\newtheorem{proposition}[theorem]{Proposition}
\newtheorem{lemma}[theorem]{Lemma}
\newtheorem{corollary}[theorem]{Corollary}
\theoremstyle{definition}
\newtheorem{definition}[theorem]{Definition}
\newtheorem{example}[theorem]{Example}
\theoremstyle{remark}
\newtheorem{remark}[theorem]{Remark}
\newcommand{\R}{\mathbb R}
\newcommand{\N}{\mathbb N}
\newcommand{\id}{\operatorname{id}}
\title{Evidence Guided Native Experiment Formation\\
\large Resource preserving strategies and a closed manufacture revision loop}
\author{}
\date{5 November 2026}
\begin{document}
\maketitle
\begin{abstract}
We study a native policy that manufactures an instrument, performs a newly
available experiment, repairs a construction-relevant characterization,
and builds a further reusable device from its actual residual stock.
The result is conditional on explicit source, access and resource promises;
it does not infer those promises from finite observations or exceed the
complete runtime's extensional closure.

For a direct-production factory with bill-of-material matrix $C$, the
exact live and advance-manufacture requirements of a contingent demand
tree are respectively $\bigvee_h Cd_h$ and $C(\bigvee_h d_h)$. They agree
for every tree exactly when every row of $C$ has at most one positive
entry. This is also exactly the confluence condition for direct
manufacture, and yields a native universal capacity normalization.
When the condition fails, a two-branch witness separates the resource
contracts. A finite deterministic specialization computes reusable-service
invariants and a transition-stable completion, proves exact commutation
with residual belief updates, and synthesizes observation-uniform
manufacturing policies. The fixed-point machinery is standard; native
execution and residual-resource obligations are retained explicitly.

A continuous construction supplies the full physical loop. Finite probes
can certify a feedback latch under independently warranted transfer-law
bounds. The assembled latch measures a delayed source relation that no
matched finite feedforward baseline can robustly recover. Its result
selects the physical wiring of a separate stateless corrector, warranted
for later trials by a fixed-source assumption. A mass-conserving family
then exhibits an exact additional-manufacture separation between repeated
interleaving and manufacture once. We separate skeptical evidence update,
conjecture retraction and observational refinement, and identify precise
prior overlap. The contribution is a constructive synthesis and boundary
analysis, not a universal theory or historical priority claim.
\end{abstract}
\tableofcontents


\section{The construction loop and the claim}\label{sec:loop-question}
A manufactured instrument can change which experiment is executable. An
outcome of that experiment can make the investigator's previous
characterization inadequate, and its replacement can determine a further
physical construction. The question is whether these steps form one
correct native execution, rather than a sequence of independently true
existence statements. In particular, the experiment must leave the stock
and state needed by the subsequent construction, and the investigator
must choose a continuation from the observations actually available.

We establish a conditional constructive account of that loop. The main
example starts with independently warranted but incompletely known leaky
cells. Four finite probes can certify a robust feedback latch, actual
connections create it, and the latch enables a delayed relational
experiment. The experiment separates two source laws whose individual
endpoint marginals agree. Its outcome selects and warrants assembly of a
separate reusable signal corrector. The corrector is a physical wiring
choice, not a name for the observation. A matched baseline cannot choose
the correct wiring uniformly without acquiring the missing relation.

The mathematical account has three linked levels. A native strategy
retains actual command witnesses and the possible residual states of all
hidden worlds compatible with the visible history. Local comparison data
compile such strategies, with a separate criterion for resource
accounting across mutually exclusive branches. A finite specialization
constructs a policy and checks its terminal reusable service. Finally,
the physical example discharges the sensing, assembly, robustness, and
continuation obligations in one model. A discrete family isolates the
advantage of repeated manufacture--experiment interleaving as the depth
increases. These are instances of the same state-threaded semantics,
not assertions that a finite transition table and a continuous plant
are identical mathematical objects.

Several ingredients are established theory. Observation-based contingent
planning, alternating and feedback refinement, contextual congruences,
and controller cost transfer have strong precedents. The finite
predecessor recursion is not claimed as a new learning algorithm. The
intended contribution is the explicit combination of native witness
extraction, branch-sensitive resource transport, operational revision,
and an evidence-dependent second construction, together with sharp
failures when one of those obligations is absent. Historical priority
for that combination is not established here.

\subsection{Four changes that must remain distinct}
Let $q$ denote the entire physical configuration: artifact and stock,
wiring, controller state, stored records, retained energy or consumables,
and accessible ports. A fixed native runtime supplies the elementary
commands, observations, and admissible composition rules. A service
contract $\mathcal S$ specifies the future inputs, allowed control,
retained state, errors, and success condition. Write
$\operatorname{Cap}_{\mathcal S}(q)$ for tasks with actual implementations
from $q$ under this contract.

A manufacturing execution $m:q_0\to q_1$ can give
\[
 T\in\operatorname{Cap}_{\mathcal S}(q_1)
       \setminus\operatorname{Cap}_{\mathcal S}(q_0).
\]
This is an installed-service gain. It need not preserve every old
capability: manufacture can consume useful stock. It is also different
from a reduction in cost, an increase in warranted knowledge, or an
increase in global runtime expressivity. The complete factory may
already implement manufacture followed by the new service.

\begin{proposition}[The whole-runtime boundary]\label{prop:whole-runtime}
Suppose a fixed native runtime contains uniform implementations of
manufacture, representations, every newly named interface operation,
observations, and retained-state passage. Suppose these implementations
preserve the stated sequencing and branching semantics. Then every
finite experiment written using the new interface compiles to a program
of the original runtime. The conclusion includes adaptive manufacture
and revision. Costs and infinite interaction transfer only under
additional cost and iteration-preservation hypotheses.
\end{proposition}
\begin{proof}
Replace each newly named operation by its native implementation,
including its initialization and actual residual state. Translate each
branch using the corresponding native observation decoder. Structural
induction on the finite experiment shows equality of boundary behavior.
Prefixing the translated experiment by the manufacture program gives an
original-runtime program. The induction proves neither a cost bound nor
an assertion about an infinite computation; each requires the extra
preservation premise stated in the proposition.
\end{proof}
This is a compilation boundary, not a no-manufacture theorem. All
subsequent strictness claims fix either the installed service regime or
the information and resources from which one common construction must
be guaranteed. They do not contradict Proposition~\ref{prop:whole-runtime}.

\subsection{What a closed loop has to certify}
A valid loop has one observation-based policy, not one successful program
per hidden world. Every possible branch retains: the actual residual
stock and physical state; the evidence and model promises that warrant
its characterization; a native continuation executable from every
compatible residual state; and a reusable final interface certificate.
Controller memory, decoding, timing, and auxiliary energy are part of
this account whenever the relevant contract charges them.

Three logical changes can accompany the physical changes. Filtering a
fixed compatible-model class strengthens its skeptical consequences.
Refuting a selected conjecture can require retracting it. Adding an
executable diagnostic can require refining an old observational
quotient. These are not the same operation. Nor is a coarsest refinement
relative to a fixed source, operation family, and diagnostics a uniquely
best choice of language or theory among arbitrary alternatives.


\section{Native execution and observation based choice}\label{sec:native-semantics}
A common semantics accommodates both a finite factory and a continuous
apparatus. Let $\Theta$ be a set of fixed model parameters and let
$P_\theta$ be the possible full physical configurations in world $\theta$.
Use the tagged state space
\[
 S=\coprod_{\theta\in\Theta}\{\theta\}\times P_\theta.
\]
The world tag records an original law or task parameter; it survives even
if two experiments merge physical states. A current information state is
a nonempty set $B\subseteq S$. Its members are possible actual residual
states, not new preparations of the original specimen.

A command $a$ is supplied by a native program and a certified terminating
execution relation
\[
 E_a(s)\subseteq\{(o,s',w):o\in O_a,\ s'\in S,\
                         w\text{ is an actual native execution }s\to s'\}.
\]
Whenever $a$ is enabled at $s$, this set is nonempty and contains every
licensed outcome, including all admitted disturbances. The endpoint
includes consumed stock and retained apparatus. A macro is admissible
only with its own termination, observation-decoding, and endpoint
certificate. Continuous bounded-duration flow, a link action, and a
finite reaction script can all be commands under this definition.
A wait for an externally arriving event is terminating only under an
explicit eventual-arrival promise; a uniform duration bound is a stronger
condition not assumed automatically.

A command is safe at $B$ only when it is enabled at every member. Its
visible outcome update is
\begin{equation}\label{eq:general-post}
 \operatorname{Post}(B,a,o)=
 \{s':\exists s\in B\ \exists w,\ (o,s',w)\in E_a(s)\}.
\end{equation}
Only nonempty sets are possible successors. Observable refusal can be a
legitimate outcome of an explicitly supplied total attempt command.
Undefinedness, an unlicensed timeout, and divergence cannot instead be
used as free observations or vacuously successful branches.

\begin{definition}[Native contingent witness]\label{def:native-policy}
A finite native contingent witness at $B$ is a rooted observation tree.
At every nonterminal node its label is one command safe throughout the
exact residual belief at that node. Its edges cover all possible decoded
outcomes. Each actual edge occurrence retains the compatible native
execution witnesses and the resulting belief from
\eqref{eq:general-post}. A leaf supplies a terminal certificate valid for
every member of its residual belief. Choices depend only on the visible
history, not on a hidden state chosen to explain that history.
\end{definition}
Finite depth can coexist with a macro whose duration has no uniform bound,
as in a promised eventually arriving second signal. A runtime with bounded
controller memory must supply enough memory to execute the selected tree
or a finite-state implementation of it. A set-theoretic tree with
uncountably many undecodable outcomes is not automatically an executable
controller.

Native execution histories themselves form a category: objects are actual
physical states and arrows are the licensed witnessed executions,
composed by concatenating at the actual common endpoint. Empty histories
are identities. Its reachability preorder forgets all arrow identity,
resource use, and observation history. A contingent witness is a branching
execution prescription over this category, not an arbitrary arrow in the
preorder. Sequential composition is grafting at compatible residual
beliefs; it is associative because actual execution concatenation and
history concatenation are associative. No choice of a fresh convenient
representative is made at a graft.

\begin{proposition}[Local native implementations lift finite strategies]
\label{prop:macro-lift}
Suppose an abstract command system has states $X$ and a relation
$R\subseteq S\times X$. For each abstract command $a$ and each available
abstract controller state, supply one native macro $M_a$, selected
without hidden-state access, and an effective decoder $d_a$ of its visible
outcome. Assume the following whenever $sRx$ and $a$ is enabled at $x$:
\begin{enumerate}
\item $M_a$ is nonblocking and terminating at $s$;
\item every native outcome $(o,s',w)$ has a matching abstract transition
      $x\xrightarrow{a,d_a(o)}x'$ with $s'Rx'$;
\item its witness $w$, actual residual resources, and any controller or
      decoding overhead satisfy the stated local implementation budget.
\end{enumerate}
Assume abstract terminal certificates are valid at every related native
state. Then any finite abstract observation-based winning policy compiles
by these macros to a native winning policy at related initial states.
It retains the concatenated native witnesses. A pathwise additive budget
transfers by summing the local bounds, and a worst-case budget is the
supremum of those sums over possible native branches.
\end{proposition}
\begin{proof}
At the root, the common macro is executable in every related hidden
state. Every native outcome decodes to a branch of the abstract policy
and leaves a related actual endpoint. Apply the same argument recursively
at that branch. Induction on tree depth reaches a valid terminal
certificate in every execution. It also constructs the native tree,
retaining the witness already taken at each step rather than choosing
another witness retrospectively. Summing the local inequalities on an
actual path proves its budget; taking the supremum proves the final
bound. The implementation's extra observations can be ignored, but they
cannot be used to choose a different abstract branch from their decoded
value.
\end{proof}
This is a sufficient refinement theorem, not an extraction algorithm for
the macros. The macros and their common observation-based selection are
real premises. Ordinary simulation with an existential implementation
chosen separately for each hidden $s$ does not supply them. The finite
completion below derives an exact special case directly from native
tables, while the resource theorem identifies a further structural
condition for moving manufacture across a branch. These additional
results prevent the general proposition from merely assuming the desired
compiler.

\begin{proposition}[Strictness needs reflection of the baseline]
\label{prop:baseline-reflection}
Let a comparison preserve the positive native policy for a task after
manufacture. If every admissible target baseline policy for that task
lifts to an admissible source baseline policy with the same relevant
resource and observation contract, a source baseline impossibility
remains an impossibility in the target. Preservation alone is insufficient.
\end{proposition}
\begin{proof}
A target baseline policy would lift to the prohibited source policy.
For failure without reflection, add a free perfect sensor or persistent
controller only to the target baseline. All source programs still
translate, even with an injective state encoding, but the target can now
solve the previously impossible task.
\end{proof}


\section{Why diagnosis and construction cannot be proved separately}
\label{sec:obstructions}
The resource-sensitive strategy requirement is essential even in the
smallest examples. Here is a counterexample in which neither uncertain
enabledness nor nondeterminism explains the failure.

\begin{proposition}[Perfect diagnosis can destroy all successful continuations]
\label{prop:destructive-diagnosis}
There is a two-world native system in which each world admits a
successful construction, a finite experiment identifies the world
exactly, and the identified world determines a successful construction
from the initial stock, yet no observation-based policy surely completes
the construction from that stock.
\end{proposition}
\begin{proof}
Let the hidden world be $\theta\in\{0,1\}$ and let the initial stock
contain one fresh specimen $r$. Initial observations are identical.
There are three commands, all enabled in either world when $r$ is fresh.
The test $e$ returns $\theta$ and consumes $r$. The command $b_i$ consumes
$r$ and makes the target exactly when $\theta=i$; otherwise it makes an
irreversible failure. No command produces the target without a fresh
specimen. The successful world-indexed construction is $b_\theta$ and
$e$ is a perfect diagnostic. Nevertheless the first informative action
of a common policy cannot be $e$, because its residual stock has no
successful continuation. If the first action is $b_i$, it fails in the
other world. Delaying does not change the available information.
Randomization cannot give sure success: every build action in its
support fails in one world, and testing still destroys the specimen.
\end{proof}
An extra charged specimen or a nondestructive test changes the conclusion.
One cannot add either silently while concatenating a diagnosis theorem
and a construction theorem. Beliefs therefore retain pairs
$(\theta,q)$, including the actual residual $q$, rather than merely the
remaining values of $\theta$.

\begin{proposition}[Worldwise witnesses are not one implementable policy]
\label{prop:uniform-choice}
Suppose two worlds share every admitted observation before the terminal
construction choice, and their sets of successful construction choices
are disjoint. No deterministic observation-based policy surely chooses a
successful construction in both worlds. The conclusion also holds for
randomized policies under a sure-success contract.
\end{proposition}
\begin{proof}
The policy receives the same observable history in both worlds and must
therefore select the same choice, or the same distribution of choices.
A sure-success choice must lie in both successful sets; a sure-success
distribution must have its support in their intersection. The
intersection is empty.
\end{proof}
This elementary quantifier obstruction will reappear as the ablation in
the physical construction. Existence of a witness at every hidden state
is weaker than existence of a uniformly executable observation-based
witness. A categorical comparison that forgets the difference cannot
preserve the relevant feasibility statement.

\subsection{The limit of finite validation}
A finite successful test cannot establish an indefinitely reliable
service for an unrestricted physical source class. For any finite test
horizon, a deterministic time-homogeneous hidden-state device can agree
with a reliable device throughout that horizon and then fail after a
larger hidden counter reaches zero. The class can consist entirely of
finite-state devices if no common state bound is supplied. Thus
stationarity alone does not remove the obstruction. The positive
continuous construction below relies on independently warranted
monotonicity, dynamics, range, and disturbance bounds. The finite
observations certify inequalities inside that class; they do not certify
the class itself. A failed sufficient certificate is reported as
inconclusive, not as proof of impossibility.


\section{When adaptive construction can be moved before observation}
\label{arc:sec}

The relevant comparison is not whether a complete adaptive program lies
outside its original runtime. It is whether the manufacture needed by that
program can be moved before a newly available observation, at the same
material budget. Individual fixed experiments can admit such a move even
when their observation-dependent assembly does not. The following class
has an exact primitive-level criterion and native witnesses.

\subsection{An explicit native factory and the timing contract}
Fix $C\in\mathbb N^{d\times m}$, with columns $c_i$, and assume
\emph{every column is nonzero}. There are raw-material species
$R_1,\ldots,R_d$ and inert product species $A_1,\ldots,A_m$. Native
factory markings are pairs $(r,a)\in\mathbb N^d\times\mathbb N^m$.
The only manufacturing rules are
\[
 \mu_i:(r+c_i,a)\longrightarrow(r,a+e_i).
\]
Delivery of one $A_i$ removes one available copy and records the public
label $i$. Products cannot be converted back, substituted for one another,
or manufactured by another route. There are no later material supplies.
All unused products and raw stock remain in the residual marking.
A rule may have any finite material multiplicities. Consequently neither
Boolean species presence nor an uncharged copy operation suffices.

An observation or demand-arrival event has an independently supplied
terminating interface. It does not consume factory stock and its law is
unchanged by unused products stored inertly. It may depend on the external
world, earlier public observations, and earlier delivered products.
Every branch treated as possible must be physically admitted by that
interface. If an observation consumes a specimen or changes construction
authority, that resource must be added to the model; the theorem does not
silently regard it as free. A previously manufactured sensor may supply
these events: its manufacture, access, retained memory and read costs are
charged in a common prefix, not supplied by the theorem.

A finite \emph{demand tree} $P$ prescribes observation-dependent finite
sequences of deliveries. For each possible complete history $h$, let
$d_h\in\mathbb N^m$ count its cumulative deliveries. This retains the
whole branch history, rather than merely the separately attainable tests.
There are two contracts, starting with raw stock $r$ and no products.
\begin{itemize}
\item In the \emph{live factory}, manufacture remains available between
observations and deliveries. A requested product can be made just in time.
\item In the \emph{frozen factory}, a finite setup is completed before the
first outcome or future demand is supplied. The factory then loses
manufacturing authority. Previously made products can still be selected,
used and delivered after observing the outcome.
\end{itemize}
The same public demand/observation tree, product services, delivery
operations and disturbance assumptions are used on both sides. The timing
of manufacturing authority is the declared difference. No setup depends
on a future outcome. Setup may depend on the known tree, unless the
stronger uniform result below is invoked. Availability is relative to
this service contract, not global runtime expressivity.

Write $\bigvee$ for coordinatewise maximum and define
\[
 D(P)=\bigvee_h d_h,\qquad
 L_C(P)=\bigvee_h Cd_h,\qquad
 F_C(P)=C D(P).
\]
These are material inventories, not elapsed-time or computational-cost
bounds.

\begin{theorem}[Exact adaptive manufacture comparison]
\label{arc:profile}
For the native factory above, a demand tree $P$ is realizable from raw
stock $r$ in the live contract if and only if $r\ge L_C(P)$, and in the
frozen contract if and only if $r\ge F_C(P)$.
Both bounds are attained by explicit finite native manufacturing and
delivery histories. Always $L_C(P)\le F_C(P)$.

The following conditions on the primitive matrix $C$ are equivalent.
\begin{enumerate}
\item Every row of $C$ has at most one positive entry.
\item The additive monotone map $C:\mathbb N^m\to\mathbb N^d$
preserves finite joins.
\item $L_C(P)=F_C(P)$ for every finite public demand tree, allowing a
binary public input with both outcomes possible.
\item Direct manufacture rewriting is confluent from every marking:
any two manufacture histories from the same marking have a common
manufacture descendant.
\end{enumerate}
The criterion is checked directly from the entries of $C$. If it fails,
two positive entries in one row produce a two-branch counterexample,
including a stock vector and an unjoinable native critical pair.
\end{theorem}
\begin{proof}
Along a history delivering $d_h$, at least $d_{h,i}$ copies of $A_i$
must have been manufactured. The only rule making them consumes $c_i$
per copy. Thus every live realization requires $r\ge Cd_h$ for every
possible $h$. Conversely, manufacture each requested copy immediately
before delivery. Every prefix consumes no more, coordinatewise, than
any possible completion of that prefix. The bound $r\ge L_C(P)$
therefore enables every requested rule along every possible history.
This is one observation-uniform program: it follows the public tree and
does not inspect the hidden world.

A frozen setup must leave an inventory $a$ with $a\ge d_h$ for every
possible $h$, since no further product can appear. Hence $a\ge D(P)$
and $r\ge Ca\ge CD(P)$. Conversely, manufacture exactly $D(P)$
before the first outcome, in any enabled order, and then follow $P$.
The stock bound enables all those manufacture steps and every delivery.
These constructions give native witnesses, not only semantic existence.
Monotonicity gives $Cd_h\le C D(P)$, proving $L_C(P)\le F_C(P)$.

If each row has at most one positive entry, that row sends a coordinate
maximum to the maximum of its values. This proves (1)$\Rightarrow$(2),
and (2) implies (3). If row $k$ contains positive entries in distinct
columns $i,j$, use two branches demanding respectively $e_i,e_j$.
Their live requirement is $r=c_i\vee c_j$, while the frozen requirement
is $c_i+c_j$. Its $k$th coordinate is strictly larger than $r_k$.
This proves (3)$\Rightarrow$(1) and also disproves join preservation.

Under (1), define
\[
 K_i(r)=\min_{k:C_{ki}>0}\left\lfloor\frac{r_k}{C_{ki}}\right\rfloor.
\]
Every column is nonzero, so this is a finite integer. Different products
use disjoint sets of raw species. Hence $CK(r)\le r$, and every maximal
manufacture history from $(r,a)$ ends at
\[
 N_C(r,a)=\bigl(r-CK(r),\ a+K(r)\bigr).
\tag{*}
\]
One can extend every manufacture history to this marking by making its
remaining copies. Termination follows because each manufacture consumes
at least one raw unit. Thus all pairs of histories have a common
descendant, proving (4). Conversely, for a row shared by $i,j$, start
from $(c_i\vee c_j,0)$. Both $\mu_i$ and $\mu_j$ are enabled, but no
descendant can contain one of each: that would consume $c_i+c_j$, which
exceeds the initial stock in the shared row. Products are never removed
by manufacture. The two successors therefore have no common descendant.
This proves (4)$\Rightarrow$(1).
\end{proof}

\begin{proposition}[One universal setup in the positive class]
\label{arc:uniform}
Under the equivalent conditions above, $N_C(r,a)$ is a finite setup
independent of any later demand tree. Put
\[
 u=r-CK(r),\qquad b=a+K(r).
\]
For all finite clients built from the stated observations and product
deliveries, the live marking $(r,a)$ and frozen marking $(u,b)$ support
exactly the same public delivery histories. In the live factory,
manufacturing preserves $(u,b)$ and delivery of $i$ decreases $b$ by
$e_i$. A delivery, with manufacture first when required, is possible
exactly when $b_i>0$. Thus this is a constructed all-client comparison,
not a relation assumed as an input.
\end{proposition}
\begin{proof}
If $\mu_i$ is enabled, disjoint row supports give
$K(r-c_i)=K(r)-e_i$. Hence its change from $(r,a)$ to
$(r-c_i,a+e_i)$ preserves both $u$ and $b$.
A delivery changes only $a_i$ and therefore decreases only $b_i$.
If $b_i>0$, either $a_i>0$ or $K_i(r)>0$; in the latter case manufacture
then delivery is enabled. If $b_i=0$, neither a stored copy nor enough raw
material for a new copy exists. The same criterion governs delivery from
the frozen inventory $b$. Observations have the same law and ignore unused
factory inventory, so induction over any finite client preserves the
public history and the displayed residual relation. This argument permits
arbitrary observation-dependent future demands and arbitrary product
multiplicities; it does not enumerate a fixed finite client library.
\end{proof}

The normalization does not preserve manufacturing work: it makes
$\sum_iK_i(r)$ items up front, whereas a live client may make fewer or
none. Zero-cost infinite manufacture is excluded by the nonzero-column
hypothesis. Infinite clients and liveness require their own semantics.
The result neither changes the runtime's extensional closure nor licenses
new physical operations, environmental promises, or indefinite storage.

\subsection{Residual resources, histories, and coherent composition}
For the exact frozen realization of $P$, after history $h$ the residual
factory marking is
\[
 (r-CD(P),\ D(P)-d_h).
\]
The just-in-time live realization has residual marking $(r-Cd_h,0)$.
These are usually \emph{different states}. Their conserved material
potentials agree:
\[
 r-Cd_h=(r-CD(P))+C(D(P)-d_h).
\tag{$\dagger$}
\]
Unused prefabricated products are recorded, not discarded or treated as
raw stock. Equation $(\dagger)$ is not a back-conversion procedure.
With a shared raw resource, those products may be unusable for a later,
unanticipated different demand. In the positive class, the universal
capacity certificate proves adequacy only for the declared delivery and
noninterfering observation clients. Adding raw-material-sensitive tests,
a new product conversion, resource-sensitive clocks, or return operations
changes the contract and requires a new comparison.

Native witnesses can be kept in the execution category whose objects are
markings and whose arrows are actual finite rule histories. If a chosen
native doctrine also has cells comparing independent firing orders, retain
those cells too. The profile map forgets information; the comparison
therefore records pairs of live and frozen witnesses, their common public
history, and their exact residual markings. It does not assert equality
of their native paths, recover forgotten allocations, or infer a section
of the witness-forgetting functor. In particular, a confluent rewrite
system has a unique normal \emph{marking}, not a unique native execution.
The earlier free-history diamond obstruction is not being contradicted.

There is a precise compositional law before profiles are minimized.
For a tree $P$ and a continuation tree $Q_h$ after each leaf $h$, write
$e_{h,k}$ for its continuation demand. Then
\[
 \begin{split}
 D(P[Q])&=\bigvee_{h,k}(d_h+e_{h,k}),\\
 L_C(P[Q])&=\bigvee_{h,k}C(d_h+e_{h,k}),\\
 F_C(P[Q])&=C\bigvee_{h,k}(d_h+e_{h,k}).
 \end{split}
\tag{$\ddagger$}
\]
These equations are exact. Three successive substitutions have the same
leaf histories and cumulative demands under the canonical reassociation
$((h,k),l)\leftrightarrow(h,(k,l))$. Native histories concatenate at their
actual residual markings, so the associativity and unit comparisons are
those of execution, and reassociation satisfies the pentagon. This is
ordinary coherent tree substitution with resource data retained, not a
claim that minimizing each component separately is strongly compositional.
Indeed one has only the safe bound
\[
 D(P[Q])\le D(P)+\bigvee_hD(Q_h)
\]
in general. Extra prefabricated inventory in this bound remains explicit.
Replacing the full family $d_h$ by its maximum before substitution can
lose the correlations needed to obtain the exact smaller answer.

For independent parallel tasks $P,Q$ with Cartesian-product possible
histories and disjoint/noninterfering service interfaces, both material
profiles are additive:
\[
 L_C(P\otimes Q)=L_C(P)+L_C(Q),\qquad
 F_C(P\otimes Q)=F_C(P)+F_C(Q).
\]
This follows from $\bigvee_{h,k}(d_h+e_k)=\bigvee_hd_h+\bigvee_ke_k$
and the same identity after applying $C$. Native realization requires
actual allocation of the summed stock to the two executions. The equality
does not supply copying, independence, concurrent throughput, or a bound
on elapsed time.

\subsection{The categorical comparison and its exact obstruction}
The ordered commutative monoids $(\mathbb N^m,+,\le)$ and
$(\mathbb N^d,+,\le)$ are thin symmetric monoidal categories. The
bill-of-material map $C$ is a strong monoidal functor. Finite joins encode
one inventory sufficient for every possible branch. Monotonicity gives
the canonical comparison
\[
 \bigvee_h C d_h\longrightarrow C\Bigl(\bigvee_h d_h\Bigr).
\]
The theorem constructs the native meaning of both endpoints and proves
that this comparison is invertible for all branch families exactly when
each raw resource is exclusive to at most one product type. It is
therefore a primitive-level criterion for an interchange of adaptive
choice and manufacture, not an assumed lifting of whole strategies.
Successive bill-of-material maps compose their monotone comparisons
coherently, but a composite can satisfy or fail the row criterion
independently of an unverified physical implementation of its factors.
Any physical interpretation still has to retain intermediates, catalysts,
resource ownership and execution witnesses.

\begin{example}[Every fixed test lifts, but the adaptive test needs $m$ times the stock]
\label{arc:gap}
Take one raw blank $R$, $m\ge2$ actuator types $A_i$, and native rules
$R\to A_i$. A manufactured sensor produces an outcome
$o\in\{1,\ldots,m\}$, each possible. A fixed test chooses $i$ before
execution, reads the sensor, and requests $A_i$. One blank suffices in
either factory: in the frozen one, manufacture $A_i$ before the read.
This compilation retains the same public read and delivered actuator,
with the actual earlier manufacturing history recorded.

The adaptive test requests $A_o$. In the live factory one blank suffices:
read $o$, then execute $R\to A_o$. In the frozen factory at least one of
every $A_i$ must already be stored, so exactly $m$ blanks and $m$
manufacturing firings are necessary and sufficient. Here
$C=(1,\ldots,1)$, $d_o=e_o$, $L_C=1$ and $F_C=m$.
With only one blank, any fixed or randomized prechoice misses an admitted
outcome; randomization cannot provide a sure guarantee. If $o$ is uniform conditional on all prechoice information and independent
of the controller's private randomization, then for deterministic
or randomized prechoice, the best success
probability is at most $1/m$; with $b\le m$ distinct preinstalled types
it is $b/m$. These probability statements use the extra stated outcome
law, whereas the sure-success material bound does not.
\end{example}


The observations in this theorem have their own noninterference contract.
It cannot be applied directly to an experiment that consumes the only
specimen or closes the factory. Such changes belong in the full residual
state of Section~\ref{sec:native-semantics}. The next finite theorem does
that explicitly. Conversely, once a sound controller supplies a finite
demand tree and a separate direct-production reserve satisfies the
hypotheses here, the displayed profiles determine exactly whether that
reserve can be committed before observation. The two results address
complementary obligations of the same contingent execution.


\section{A computable completion for adaptive manufacture}
\label{msc:sec:completion}

There are three different outputs to construct. A response quotient says
which source states an interface can identify. A learning policy says what
can be established on the available specimen. A manufacturing policy says
what can still be built from its actual residual state. The last two do not
follow from the first. This section gives a finite algorithm that produces
all three together, or a finite obstruction, starting with native process
tables rather than an assumed correct quotient.

The positive result is deliberately relative to a supplied finite native
model class, finite stock bound, and declared service boundary. It constructs
neither the physical laws nor an unrestricted new language from finite data.
The completion is canonical only after these inputs and the retained old
characterization are fixed. Its substantive conclusion is the commuting
comparison between completed source states and resource-sensitive adaptive
manufacturing policies, together with extracted native executions.

\subsection{Independent inputs: worlds, residual stock, and primitive commands}
Let $\Theta$ be a finite set of possible, stationary native worlds. For each
$\theta$, let $P_\theta$ be a finite set of physical configurations. A
configuration includes installed wiring, remaining stock, damaged and spent
parts, retained registers, and any relevant control mode. Put
\[
 S=\coprod_{\theta\in\Theta}\{\theta\}\times P_\theta.
\]
World tags are never discarded merely because two physical configurations
look alike. Initial evidence supplies a nonempty belief $B_0\subseteq S$,
not a distinguished true world. The factory controller has supplied capacity to store a subset of $S$
and the finite budget counter used below; this storage and computation are
explicit parts of the access contract. No stronger bounded-memory theorem
is claimed. Device registers and any smaller physical controller bound
that restricts an action are included in the source configuration.

A finite command alphabet $A=A_{\rm fac}\sqcup A_{\rm svc}$ has supplied
partial deterministic transition tables
\[
 T_a:S\rightharpoonup O_a\times S.
\]
Each command is physically executable and terminates on its declared
domain. A factory command can link, unlink, actuate, probe, expend a specimen,
or seal a device. It has a positive integer charge $c(a)$. Physical stock
changes are recorded by $T_a$, in addition to this scalar accounting.
Transitions preserve the world tag. A command not enabled at a possible
state may not be attempted safely. Its disabledness is not automatically
an observable outcome. A genuinely available failure-reporting command
must instead have an explicit transition and output.

The tables can themselves be generated from a bounded collection of native
components and their gate, reaction, or finite-state laws. Supplying candidate
tables is substantially weaker than supplying the true table, but it is
still a source-model promise. Exact observations eliminate candidates inside
that promise; they do not establish its unrestricted physical validity.

For a uniformly enabled command define the actual residual belief
\[
 \operatorname{Post}(B,a,o)=
 \{s':\exists s\in B,\ T_a(s)=(o,s')\}.
\]
Only nonempty successors are possible observations. This update preserves
all world--configuration correlations; replacing it by a posterior over
worlds with fresh stock would change the problem.

\subsection{Generate the reusable services from the native tables}
Fix finite boundary alphabets $U,V$, with $U\ne\varnothing$, and a specified
service mode $S_{\rm svc}\subseteq S$. For each $u\in U$, the primitive service
command $a_u$ supplies the input and returns a symbol in $V$. Its transition
may change hidden device state. These service labels may be environment
stimuli rather than controller-visible commands: only $A_{\rm fac}$ is
chosen by the factory policy. In that case $T_{a_u}$ is the actual composite
of the environmental source and installed device; checking it for every $u$
does not give the device direct access to an otherwise hidden stimulus.
The service controller does not manufacture
an uncharged replacement, reset the world, or obtain a future input in
advance. A more elaborate service protocol needs an explicit finite-control
implementation and a checked request-to-response transition table; the
present construction does not obtain such a macro from an arbitrary
behavioral oracle.

For every function $f:U\to V$, compute
\begin{align*}
 I_f^{(0)}&=S_{\rm svc},\\
 I_f^{(j+1)}&=\{s\in I_f^{(j)}:\ \forall u\in U,\
       T_{a_u}(s)=(f(u),s')\text{ for some }s'\in I_f^{(j)}\},\\
 I_f&=\bigcap_{j\ge0}I_f^{(j)}.
\end{align*}
At most $|S|$ strict deletion rounds occur for each $f$. There are only
$|V|^{|U|}$ functions to test. Thus this computes the entire finite family of
pure reusable services at the declared boundary, not just one desired
operation or a supplied service oracle. Empty $I_f$ are discarded.

\begin{lemma}[Exact reusable-service certificates]\label{msc:lem:service}
For an initial state $s\in S_{\rm svc}$, $s$ belongs to $I_f$ if and only if every finite sequence of service
inputs is executable and returns the pointwise values of $f$, remaining in
service mode after every request. The finite invariant $I_f$ is an
independently checkable certificate. Distinct functions have disjoint
nonempty invariants.
\end{lemma}
\begin{proof}
The fixed-point equation proves the forward direction by induction on the
number of requests. Conversely, the states satisfying the claimed universal
property form a post-fixed point of the deletion operator. Induction on $j$
puts them in every $I_f^{(j)}$, hence in $I_f$. If $f\ne g$, choose $u$ with
$f(u)\ne g(u)$. Determinism prevents a common state from satisfying both
one-request requirements.
\end{proof}

Write $\gamma(s)=f$ when $s\in I_f$, and $\gamma(s)=\bot$ otherwise. A
manufacturing target may be one discovered $f$, any specified subset of the
discovered services, or the complete list. The boundary and service mode
are supplied, but the functions actually realized and their invariant
representations are computed from native transitions.

\subsection{Complete the old characterization before using it as a model}
Let $q^0:S\to Q^0$ be the retained old characterization. Include in its seed
any exact public resource, port, or control labels that must be retained.
It need not already support the new fabrication or probe commands. Start
with the partition
\[
 sE_0t\quad\Longleftrightarrow\quad
     q^0(s)=q^0(t)\ \text{ and }\ \gamma(s)=\gamma(t).
\]
For a partition $E$, give each state the finite signature
\[
 \sigma_E(s)=\left([s]_E,
   \left(\begin{cases}
      \mathsf{disabled},&T_a(s)\text{ undefined},\\
      (o,[s']_E),&T_a(s)=(o,s')
   \end{cases}\right)_{a\in A}\right).
\]
Set $E_{j+1}=\ker\sigma_{E_j}$ and stop when no split occurs. Let $E_*$ be
the result, $Q=S/E_*$ and $q:S\twoheadrightarrow Q$.

\begin{proposition}[Finite manufacture-stable completion]
\label{msc:prop:completion}
The construction terminates after at most $|S|-|S/E_0|$ strict rounds. It
returns the coarsest quotient refining $E_0$ for which every primitive
command has well-defined enabledness, output and successor class. In
particular there are induced partial tables
\[
 \overline T_a:Q\rightharpoonup O_a\times Q,
 \qquad \overline T_a(qs)=(o,qs')\quad\text{when }T_a(s)=(o,s'),
\]
and a unique map $\rho:Q\to Q^0$ with $q^0=\rho q$.
Every strict split has a finite certificate: a primitive-command context
ending in a seed difference, output difference, or enabledness difference.
Its depth is at most the refinement round. Such a table certificate is an
executable separating experiment only if its prefix and final observation
are jointly admissible on the compared residual states.
\end{proposition}
\begin{proof}
Each strict round increases the number of blocks. At stabilization the
signature equation gives exactly the stated descent conditions. If $F$ is
any equivalence refining $E_0$ with these conditions, induction gives
$F\subseteq E_j$ for every $j$: $F$-related states have the same outputs,
enabledness and $F$-related successors, hence the same next signature.
Therefore $F\subseteq E_*$. This proves coarseness without assuming a
desired quotient. Store the responsible command whenever two states split;
following its successors to a previous split, or to a directly differing
output/guard/seed, gives the certificate and depth bound. A disabledness
entry is a model fact, not permission to perform an unsafe action.
\end{proof}

This is deterministic partition refinement, not a new minimization theorem.
Its role is to derive the completed interface from the source. In particular,
closure only under the previously installed readouts is insufficient: the
signature must include native manufacturing transitions whose successors
can admit new readouts. The more substantial issue is whether this quotient
can be used by a single experimental history, as follows.

\subsection{The comparison with adaptive manufacture is exact}
For a nonempty raw belief $B$, put $q_!B=\{q(s):s\in B\}$. Define quotient
belief updates using $\overline T_a$ in the same way as raw updates. A
quotient controller is a finite observation-based program whose instructions
are the original native command labels, not arbitrary arrows inferred from
endpoint reachability.

\begin{theorem}[Completion, evidence update and native realization]
\label{msc:thm:joint}
The preceding algorithms compute $Q$, its tables, every reusable-service
invariant, and the following exact comparison.
\begin{enumerate}
\item A factory command is enabled throughout $B$ if and only if it is
enabled throughout $q_!B$. Its possible output set is the same on both
sides, and, for each output,
\[
 q_!\operatorname{Post}(B,a,o)
     =\operatorname{Post}(q_!B,a,o).                 \tag{MC}
\]
Moreover $B\subseteq I_f$ if and only if $q_!B\subseteq q(I_f)$.
\item For every integer budget $b$, a single observation-based native
factory policy can take $B$ to residual states certified for $f$ with
worst-case charge at most $b$ if and only if such a policy exists from
$q_!B$ in the computed quotient. The same command-and-branch program works
on both sides, and every leaf has a native endpoint in $I_f$. Hence the
minimum worst-case manufacture cost is preserved and reflected.
\item These properties are decidable. A successful computation returns a
finite policy tree, every posterior residual belief, and every terminal
invariant certificate. A failed bounded computation returns a finite
counterstrategy certificate. The successful policy, followed by the
service commands, realizes $f$ for every future finite input sequence,
with no fresh specimen or hidden reset between the experimental and
manufacturing stages.
\end{enumerate}
\end{theorem}
\begin{proof}
For (1), all members of one $E_*$-class have the same domain flag and
output for a fixed action, and their successors lie in the same class.
This gives enabledness and output equality. Each element on the left of
(MC) has a raw predecessor and hence belongs to the right. Conversely, a
quotient predecessor class on the right has a representative in $B$;
transition stability sends that representative to the stated successor
class with the stated output. Thus it also lies on the left. The invariant
claim holds because $\gamma$ is a seed, so $I_f$ is saturated under $q$.

For (2), induct on a finite policy tree. At a leaf the invariant claim is
(1). At an internal node, enabledness and possible outputs agree, and (MC)
identifies every updated belief; apply induction to each continuation.
Command labels and charges are unchanged, so so is worst-case cost. This
also applies to a terminating observation-based policy not initially
presented as a finite tree: finite branching and termination on every
possible run imply a finite reachable tree by K\"onig's lemma. A bounded
cost policy is already depth-bounded because charges are positive.

For a direct algorithm, let $\mathcal B$ be all nonempty subsets of $Q$ and
write $G_f=\{C\in\mathcal B:C\subseteq q(I_f)\}$. Compute
\begin{align*}
 W_{f,0}&=G_f,\\
 W_{f,b}&=G_f\ \cup\ \{C:\exists a\in A_{\rm fac},\
  a\text{ enabled throughout }C,\ c(a)\le b,\\[-2mm]
 &\hspace{38mm}\forall o\text{ possible},\
  \operatorname{Post}(C,a,o)\in W_{f,b-c(a)}\}.
\end{align*}
Positive charges make this a finite recursion. Induction on $b$ proves its
exactness; storing a witnessing command extracts the successful tree.
For a losing $(C,b)$ and each uniformly enabled affordable command, retain
one possible successor outside the appropriate earlier winning set. Along
with the nonterminal verdict $C\notin G_f$, these finite entries certify
that every proposed first move admits a failing continuation. Expanding the
same tree in raw tables retains actual endpoint witnesses in every world.
Finally Lemma~\ref{msc:lem:service} supplies arbitrary later reuse from those
very endpoints, rather than from a pre-experiment stock.
\end{proof}

There are at most $2^{|Q|}-1$ belief states. This is an effective theorem,
not an efficiency claim. Unbounded sure termination is also decidable by
the least reachability attractor on this finite belief space. Every strict
attractor round adds a belief; a winning rank decreases after every action.
Thus a winning strategy uses at most $2^{|Q|}-1$ actions per run, and at
most $c_{\max}(2^{|Q|}-1)$ charge. Searching bounded costs through that value
also computes the optimum over all surely terminating policies. A belief can be stored using $|S|$ bits in the raw representation, plus
finite command and budget registers. These resources must be supplied.
A stricter controller-memory synthesis problem needs its own game and
is not solved merely by putting a memory number in a state label.

\paragraph{What has actually been learned.}
After transcript $h$, the interface state is the observed posterior
$q_!B_h$, not a guessed representative of $Q$. If this set is a singleton,
the completed current-state code is known, and the quotient tables directly
predict all later primitive transitions and support ordinary shortest-path
construction. If it is not a singleton, the extracted controller is still
uniform over all remaining codes. It may realize the service without
identifying the entire world. Neither conclusion reconstructs distinctions
that a destructive probe erased from the physical stock.

\paragraph{Coherence without erasing execution histories.}
Equation (MC) is preserved under sequencing and observation branching, so
its comparison extends to every finite native program by structural
induction. Concatenated certificates use the actual residual belief of the
first program as the input of the second. At the process level, retain the
command-labelled syntax and its execution traces. This gives a compositional
compiler for those programs; it does not choose a functorial native history
for each arrow of a reachability preorder. Distinct paths through a diamond
remain distinct witnesses even if their endpoints have the same code.
Independent parallel composition requires a product runtime with genuinely
separate resources. Shared parts belong in a joint source configuration.


\paragraph{Relation to the general native semantics.}
Here $E_a$ has the unique output and endpoint supplied by $T_a$, together
with its native command witness. The quotient comparison is derived
from those tables; it is not an arbitrary semantic operation whose
compiler is assumed. The computed quotient is a model of possible
configurations, not an observed value of the unknown system. Actual
experimental learning consists in updating $B$ and hence $q_!B$.

This is an exact finite deterministic specialization. A robust invariant
region of a continuous plant supplies a sound finite interface, but does
not alone reflect every continuous manufacturing strategy or its optimal
cost. The following physical example uses the general relational
implementation proposition and an independent analytic baseline
obstruction; no finite-state completeness theorem is silently applied to
all its continuous dynamics.



\section{A manufactured memory revises a source law and builds a second device}
\label{arc:sec:latch-corrector}

The next construction closes a loop left open by the delayed-comparison
example in \emph{Endogenous Capability Growth}. It does not invent another
memory mechanism. It uses the native, robust latch constructed there to
obtain new evidence, determines exactly which earlier characterization is
no longer sufficient, and uses the resulting information to manufacture
another reusable device. The new device is physically elementary: its
nontrivial achievement is choosing the \emph{correct} wiring from an
unknown source law, not adding either polarity function to the global
runtime's extensional repertoire.

\subsection{Independent source, access, and wiring promises}

In this section $\Theta$ denotes only the source parameter family; any
unknown cell-transfer laws and continuous residual states are separately
retained in the full native information state. Let
$\Theta=\{-1,+1\}$. A source has one unknown parameter $\theta\in\Theta$
which remains fixed throughout calibration and all subsequent service
trials covered by the guarantee. On trial $j$ it supplies an early bit
$b_j\in\Theta$, and, after a previously unspecified finite delay, a late bit
\begin{equation}\label{arc:eq:source-law}
                 c_j=\theta b_j.
\end{equation}
The calibration and service guarantees below hold for every sequence of
bits, not only random sequences. To make the particular old statistical
characterization explicit, suppose additionally that each $b_j$ is a fresh
fair bit. Then either source has the same one-time laws
\[
 \mathcal L(b_j)=\mathcal L(c_j)=\tfrac12\delta_{-1}+\tfrac12\delta_{+1}.
\]
The old characterization $q^0:\Theta\to\{\mathsf{fair}\}$ records those
one-time laws. It is adequate for those questions and is not asserted to be
adequate for joint temporal observations, fabrication, or the full runtime.
The fair-bit promise is needed only for this characterization comparison.

The following access and manufacture contract is explicit.
\begin{enumerate}
\item A native manufacturing word $m$ assembles the robust latch proved in
\emph{Endogenous Capability Growth}, before the calibration bit arrives.
Its certified regions $B_{-1},B_{+1}$, finite write procedure, common hold
command, and positive-margin read satisfy that manuscript's reusable-bit
contract. The permitted disturbance and read errors are those of the
certificate, uniformly for every finite hold duration.
\item The early bit drives the finite write procedure. Thereafter all
external bit-dependent controller state and every external copy of the
early bit are removed or reset. Holding uses the same command for both
bits. The late bit arrives only after the write procedure has completed.
The hardware and its positive read tolerance are fixed independently of
the delay. Source bits arrive on reliable, typed digital ports.
\item A native two-rail connector is available for subsequent fabrication.
Bits have a fixed encoding $+1=(1,0)$ and $-1=(0,1)$ at fixed, labelled
boundary rails. Installing straight connections realizes $C_{+1}(c)=c$;
installing crossed connections realizes $C_{-1}(c)=-c$. These are actual
connections at the same labelled output ports, not a change in the
interpretation of an unchanged output. Their native wiring and operating
correctness are independently supplied. The latch theorem alone does not
prove the electrical correctness of an arbitrary connector.
\item After receiving the late bit, a stateless calibrated decision and
fabrication controller may read the latch, compute a Boolean product, and
install one of these two fixed connector patterns. Its selection is based
only on currently accessible observations and fixed prior information.
The selected wiring can subsequently retain its setting. No bit-dependent
structural change is allowed from the first arrival of the early bit until
the arrival of the late bit, including the finite write phase. Otherwise
connector wiring would itself be an uncharged way to remember that bit.
\end{enumerate}
The model includes the latch, the connector stock, fabrication actions,
write and read operations, source access, and the operating power supply
as explicit resources. No numeric total energy bound is asserted. It does
not claim that indefinite mathematical holding has zero energy cost or
unlimited physical lifetime.

\begin{proposition}[An explicit native certificate for the required latch]
\label{arc:prop:latch-certificate}
Take two leaky cells, fixed gain-$2$ memoryless input stages, public write
selectors, two latch link assemblies, and a disjoint reserved two-rail
corrector blank with two conductors. Under hold, initially grounded
inputs give $\dot x=-x+\eta_x$ and $\dot y=-y+\eta_y$. Execute the two native
link actions connecting the output of each cell to the other's hold-input
gain stage. The resulting hold equations are
\[
 \dot x=-x+\operatorname{sat}(2y)+\eta_x,
 \qquad
 \dot y=-y+\operatorname{sat}(2x)+\eta_y,
\]
where $\operatorname{sat}(u)=\max(-1,\min(1,u))$ and
$|\eta_x|,|\eta_y|\leq\varepsilon<1/2$. Put $M=1+\varepsilon$ and supply
initial states in $[-M,M]^2$. A write of sign $b$ selects the constant drive
$b$ at both cells for
\[
                  T\geq\log\frac{4}{1-2\varepsilon}.
\]
Then the hypotheses required of the latch above hold with
$B_+=[1/2,M]^2$, $B_-=[-M,-1/2]^2$, and any read-error bound $\nu<1/2$.
For example $\varepsilon=1/4$ and $T=\log 8$ are valid fixed parameters.
\end{proposition}
\begin{proof}
The square $[-M,M]^2$ is invariant for both writes and hold: the derivative
at either upper face is at most $-M+1+\varepsilon=0$, with the reversed
inequality at a lower face. Globally Lipschitz state laws and bounded
measurable disturbances give unique global absolutely continuous
solutions. During a positive write, variation of constants gives
\[
          x(t),y(t)\geq1-\varepsilon-2e^{-t}.
\]
The displayed duration makes the right side at least $1/2$. On $B_+$ both
cross-coupled inputs saturate to $1$; at an inner face the derivative is
at least $1/2-\varepsilon>0$, and at an outer face it is at most zero.
The rectangle is therefore invariant under every admitted disturbance.
Negative symmetry proves the other write and box. A noisy read of either
coordinate has the correct sign when $\nu<1/2$.

The links do not initialize a bit or provide an extra hidden state. The
first link makes a feedforward chain; the second closes a dynamically
guarded feedback cycle. The construction is exactly the two-cell
certificate from \emph{Endogenous Capability Growth}, reproduced here
to exhibit the physical premise rather than postulate a native memory
token.
\end{proof}

\subsection{Make the first manufacture conditional on actual finite evidence}
The preceding explicit transfer law is one way to warrant the latch. The
stronger experimental version does not assume that transfer law is known.
For each of the same two cells assume independently
\[
 \dot x_i=-x_i+\phi_i(u_i)+\eta_i,\qquad
 |\phi_i(u)|\leq1,\qquad |\eta_i(t)|\leq\varepsilon,
\]
where $\phi_i$ is unknown, time independent, nondecreasing and Lipschitz.
Set $M=1+\varepsilon$ and assume the actual cells initially lie in
$[-M,M]$. A clamp can hold either cell's input at a specified value, and a
calibration read has error at most $\nu_{\rm cal}$. Probing and subsequent
linking do not alter the promised transfer laws or cause hidden damage.
These structural and access promises are not inferred from the samples.

\begin{proposition}[Four probes warrant the first manufactured instrument]
\label{prop:four-probe}
Choose $\beta\in(0,1)$ and a settling time $\tau>0$. For each cell clamp
its input successively at $\beta$ and $-\beta$ for $\tau$, and retain the
endpoint reads $r_i^+,r_i^-$. Put
\[
 \Delta=2e^{-\tau}+\varepsilon+\nu_{\rm cal},\qquad
 c_* =\min_{i=1,2}\{r_i^+-\Delta,-r_i^--\Delta\}.
\]
If $c_*>\beta+\varepsilon$, executing the two cross-links on these same
cells yields robust invariant regions
\[
 [\beta,M]^2\quad\text{and}\quad[-M,-\beta]^2.
\]
A common-sign clamp at $\pm\beta$ for
\[
 T\geq\log\frac{1+c_*}{c_*-\varepsilon-\beta}
\]
writes either region from the actual post-probe state. Any subsequent
current-coordinate read error strictly below $\beta$ decodes the bit for
every finite hold duration. Failure of the strict inequality returns an
inconclusive certificate, not a negative verdict about all possible
latches.
\end{proposition}
\begin{proof}
The ambient interval is invariant. Under a constant input, variation of
constants and the ambient bound give
\[
 |x_i(\tau)-\phi_i(u)|
 \leq (M+1)e^{-\tau}+\varepsilon(1-e^{-\tau})
 =2e^{-\tau}+\varepsilon.
\]
Consequently $\phi_i(\beta)\geq c_*$ and
$\phi_i(-\beta)\leq-c_*$. On a lower positive-box face, cross-coupling
and monotonicity give
$\dot x_i\geq-\beta+c_*-\varepsilon>0$; the range bound gives inward
motion at the outer faces. The negative-box inequalities are reversed;
oddness is not needed. During the positive clamp,
\[
 x_i(T)\geq c_*-\varepsilon-(1+c_*)e^{-T}\geq\beta,
\]
and the negative clamp has the analogous upper bound. These bounds hold
from every ambient state, including the actual endpoint of the fourth
probe. Reading with error below $\beta$ distinguishes the boxes.
\end{proof}
The read records have finite encodings, with rounding included in
$\nu_{\rm cal}$; a certified rational upper bound for $\Delta$ may replace
its exact value. Thus the strict acceptance comparison is executable.
The native policy first performs the four probes, tests the recorded
inequality, and then executes the cross-links if certified. It retains
only a static setup certificate before the future calibration bit
arrives. Neither the probes nor the links reset a cell. The proof
constructs the write operation from whatever states they actually leave.
The final corrector blank and its two conductors are reserved throughout
these operations, physically disjoint from the two latch link assemblies.

For a guaranteed-success concrete subclass, take the actual transfer
$\phi_i(u)=\operatorname{sat}(2u)$, but do not reveal it to the certificate
algorithm. Use $\beta=1/2$, $\varepsilon=1/10$,
$\nu_{\rm cal}=1/40$, and $\tau=5$. The conservative rational bound
$\overline\Delta=3/20$ exceeds the exact $\Delta$. Every permitted record
then gives $c_*\geq7/10>3/5$, using $\overline\Delta$ in the test. A write
time $T=3$ is safe, since its worst certified ratio is $17<e^3$; read
error $1/4$ is admissible. Thus the complete first manufacture succeeds
for every permitted calibration disturbance in this subclass, not merely
for favorable records. For the wider monotone class the same policy is
sound when it accepts and may return inconclusive.

\begin{theorem}[Manufacture, warranted revision, and a reusable corrector]
\label{arc:thm:corrector}
Under the preceding promises, including a successful latch certificate,
there is a native observation-uniform policy
with the following properties.
\begin{enumerate}
\item It manufactures the latch before calibration, uses exactly one
early--late source trial, and returns $d=\theta$ for every admitted finite
delay, bit, process disturbance, and read error.
\item Its evidence refines the source-parameter family from $\Theta$ to the
singleton $\{d\}$. The coarsest refinement of $q^0$ sufficient for the
new diagnostic is the two-class characterization $q^*(\theta)=\theta$.
Both old marginal-law claims are retained.
\item It then manufactures the connector $C_d$. Without consulting the calibration bit or using the latch in subsequent
service,
the connector returns the earlier source bit on every future trial:
\[
                         C_d(c_j)=b_j.
\]
It can be reused for any finite number of future trials, subject to the
same fixed-source and connector operating promises.
\item Neither the old characterization alone nor an uninformative
calibration transcript permits a uniformly warranted choice of one of
these two correctors for the entire family $\Theta$.
\end{enumerate}
\end{theorem}
\begin{proof}
Manufacture $m$, write the early calibration bit $b$, remove the external
bit-dependent controller state, and apply the common hold command. At
arrival of $c$, the certified latch read gives $r=b$ for every admitted
delay, disturbance, and observation error. Therefore the stateless product
\[
                         d=rc=bc=b(\theta b)=\theta
\]
is exact. The finite write and read programs and the manufacturing word
are the native witnesses; the invariant-region theorem, not a numerical
trajectory or a finite collection of successful holds, supplies the
arbitrary-delay guarantee.

The source promise says that an observed $d$ is compatible with precisely
the parameter $\theta=d$. Hence exact filtering gives the singleton stated
in (ii). The new diagnostic separates the two points of $\Theta$, so any
quotient carrying it must separate those points. The identity quotient
does so and still maps to the old constant quotient. This proves the
coarsest-refinement assertion. Refinement does not refute either old
one-time marginal law; it refutes only their proposed sufficiency for the
enlarged question family.

The observed value $d$ selects a physical straight or crossed connection.
For every future trial, stationarity of $\theta$ and the native connector
law give
\[
                C_d(c_j)=d c_j=\theta(\theta b_j)=b_j.
\]
The connector is stateless and its wiring is unchanged by use. Thus finite
reuse is ordinary composition of its native executions, with no need to
copy or replay the calibration specimen. The latch can be retained as a
separate resource, but subsequent correctness does not depend on it. The two latch
links and both cells remain installed; the separate corrector blank and
its two reserved conductors become the selected fixed connector. No
calibration command consumes that reserve. Source and clamp accesses,
controller erasure, waiting energy and fabrication work are real costs
under their supplied operating contracts; no scalar optimum for them is
claimed.

Finally, $C_g$ is a correct normalizer for source $\theta$ exactly when
$g\theta=1$, equivalently $g=\theta$. The valid implementation sets are
$I_{+1}=\{C_{+1}\}$ and $I_{-1}=\{C_{-1}\}$, with empty intersection.
A selector seeing only a common old characterization or common transcript
must make the same selection in both worlds, and therefore fails in one.
This is the evidence-fibre obstruction applied to an actually fabricated
follow-on device, rather than only to a reported model label.
\end{proof}

The entire policy is observation-uniform. Before the diagnostic it uses
the same manufacturing, write, and hold procedures in both possible
worlds. Only a public outcome of an executable test selects the final
connection. The early bit is not retained by a separate experimental
controller. The permitted persistent objects are the manufactured latch
while calibrating and the subsequently manufactured connector setting.

\paragraph{Refining knowledge is not retracting a warranted old law.}
The old marginal characterization is correct in both worlds. If a
researcher additionally proposes the defeasible conjecture $c=b$, a
diagnostic result $d=-1$ refutes that conjecture and supports $c=-b$
\emph{within the two-law source promise}. These are different changes:
the skeptical compatible family is narrowed, while a selected,
previously unwarranted conjecture is retracted. A result $d=+1$ does not
establish stationarity beyond the supplied source class.

There is an explicit instance of the earlier revision-span construction.
For a field $k$ of characteristic different from two let
\[
 F=k[b,c],\quad
 J_+=(b^2-1,c-b),\quad J_-=(b^2-1,c+b),\quad
 K=(b^2-1,c^2-1).
\]
Then $J_+\cap J_-=K$, and refuting the selected repeat theory yields
\[
                 F/J_+\ \longleftarrow\ F/K\ \longrightarrow\ F/J_-.
\]
Indeed, modulo $K$, evaluation on the four bit pairs identifies $F/K$
with $k^4$; $J_+/K$ and $J_-/K$ vanish on complementary pairs, and have
zero intersection. The common fragment retains the bit equations while
discarding the incompatible selected phase law. This is an application
of the existing common-ideal theorem, not a new revision theorem. The
native latch test supplies the experimentally warranted choice of side
that the algebraic span itself cannot supply.

\begin{theorem}[A matched robust baseline cannot make the correct choice]
\label{arc:thm:corrector-obstruction}
Replace the latch in the preceding one-trial calibration contract by any
fixed finite feedforward architecture satisfying the fading hypotheses
of \emph{Endogenous Capability Growth}, with a fixed positive
current-state read-error tolerance. During holding admit no external copy
of the early bit and no bit-dependent persistent controller. Require all
architecture and connector settings to stay fixed between the arrival
of the early and late bits, including the write phase. After the late bit
arrives, permit
connector fabrication to depend only on that bit, the admitted noisy
current-state read, and fixed prior information.

No such calibration-and-fabrication procedure can always manufacture the
correct $C_\theta$ for both fixed source laws, all calibration bits, and
every finite delay. In particular, the positive result is a strict gain
for this matched calibration-to-construction task, even though both
connector functions were already individually constructible.
\end{theorem}
\begin{proof}
If such a procedure existed, apply it to an arbitrary early--late pair
$(b,c)\in\Theta^2$ and regard that pair as a permitted calibration trial
of the fixed source law $\theta=bc$. The guarantee would produce
$C_{bc}$. Evaluating its installed service at $+1$ returns $bc$, and
therefore decides whether $b=c$.

This would solve exactly the arbitrary-delay robust comparison task
excluded for that whole feedforward baseline by the earlier fading
theorem. The post-arrival fabrication does not escape the obstruction:
for sufficiently long delay the two relevant current-state read-error
sets intersect. Give the procedure the same late bit and the same
admissible current-state observation in both cases. Its available data
and hence its selected connector are then identical, although the two
cases require different connections. A randomized procedure required
to be surely correct has the same contradiction for its output support.
\end{proof}

For completeness, the relevant fading argument does not depend on one
particular matrix. Topologically order the finitely many dynamic cells.
With common hold inputs and zero disturbances, a source cell's state
difference is a constant times $e^{-t}$. If $z_i$ denotes an absolute
state difference at a later cell, fixed Lipschitz interconnections give
\[
 z_i(t)\leq e^{-t}z_i(0)+
   \int_0^t e^{-(t-s)}\sum_{j<i}L_{ij}z_j(s)\,ds.
\]
Induction bounds every $z_i(t)$ by a finite polynomial times $e^{-t}$,
so the entire state difference tends to zero. Positive-radius read-error
sets consequently intersect at a sufficiently late common read time.
The late bit can be chosen $+1$ in both cases, with early bits $+1$ and
$-1$, respectively. These are legal calibration trials of $\theta=+1$
and $\theta=-1$, and require different correctors. This proves the
indistinguishability used above directly from the native baseline laws.

The obstruction concerns the stated one-trial protocol and the stated
whole feedforward baseline. It does not exclude another prior instrument,
an externally recorded early bit, a persistent adaptive controller,
bit-dependent rewiring during writing or holding, exact unbounded-precision
sensing, or a differently warranted multi-trial protocol. Those change the
resource or access contract. In particular, the target is not a global
expressivity separation from the runtime that can manufacture the latch.

\begin{proposition}[Why stationarity is not learned from the calibration]
\label{arc:prop:stationarity}
If the source class instead permits an arbitrary phase sequence
$c_j=\theta_j b_j$, no finite calibration transcript alone warrants a
single installed polarity corrector for all future trials. The same
failure holds after any finite number of perfectly observed calibration
trials, even when every calibration phase is measured without error.
\end{proposition}
\begin{proof}
Choose two source worlds with identical bits and phases on all completed
calibration trials. On the next service trial choose equal $b$ and
opposite phases. Their calibration transcripts and installed connector
choices agree, but the required correct polarity differs. Thus one of
the two worlds violates the promised service. The fixed-$\theta$ premise
in Theorem~\ref{arc:thm:corrector} rules out precisely this pair of worlds;
it is a source warrant, not a fact inferred from finitely many trials.
\end{proof}

Similarly, allowing a nonzero probability that $c_j\ne\theta b_j$ replaces
the exact filtering claim by a statistical inference problem. One trial
then need not identify $\theta$, and connector reliability needs an
explicit error model, access assumptions, confidence criterion, and
service guarantee. The robust latch certificate protects storage and
readout; it does not certify the source law or eliminate source noise.

\subsection{What is transported between the stages}

Let a knowledge--resource interface record a compatible source family
$B\subseteq\Theta$, the installed artifacts and stock, and the actual
residual access state. The construction has the typed form
\[
 (\Theta,\text{stock})
 \xrightarrow{m}
 (\Theta,\text{latch installed})
 \xrightarrow{\text{write--hold--read and }d=rc}
 \coprod_{d\in\Theta}(\{d\},\text{calibrated latch})
 \xrightarrow{(\kappa_d)_d}
 \coprod_{d\in\Theta}(\{d\},C_d).
\]
Here the coproduct records an actual observed outcome and its compatible
branch, not permission to choose the branch containing the desired
world. The transition $m$ changes installed access without changing the
compatible family. The diagnostic changes that family. The final native
transition consumes the declared connector stock and preserves the
evidence required for its correctness claim. Their composition retains
the manufacturing word, the calibrated native test, its actual output,
the selected wiring word, and all residual resources.

This representation can live in a category of certified finite native
strategies, or in a contextad/module model after its endpoint transport,
pure-map, multiplication, and coherence premises have been independently
verified. The diagram alone does not furnish those premises. Under a
change of doctrine, a certified compiler must preserve the latch's
representation regions and read guarantee, the source relation, the
fixed labelled connector ports, actual observation branches, resource
transitions, and the native composition witnesses. Those conditions
transport the positive policy. To transport the strict impossibility
claim one also needs reflection of the entire target baseline with its
noise and memory restrictions.

\paragraph{Exact dependency and residual.}
The fabrication and robustness of the latch, feedforward fading, and the
delayed-comparison separation are prior results of
\emph{Endogenous Capability Growth}. Exact compatible-world filtering and
evidence-fibre selection are prior results of the diagnosis and revision
notes. The two-point quotient calculation and the connector algebra are
elementary. The result added here is the complete, witnessed implication
from a manufactured instrument, through an experimentally forced
characterization refinement, to a uniformly warranted second physical
artifact, together with the reduction proving that the matched earlier
architecture cannot carry out that construction task. No historical
priority or new general minimization theorem is claimed.


% theorem/proposition/corollary/definition/remark environments.
\section{A native separation between interleaving and manufacture once}
\label{ing:section}

The following finite family combines actual construction, experimental
acquisition, revision, and another construction in one policy. Its claim
is about a temporal restriction and charged manufacture in a fixed native
substrate. It is not a claim that the complete policy exceeds the runtime's
extensional closure, invents its own elementary physics, or discovers an
unanticipated explanatory language.

\subsection{The target and the unknown world}
Fix integers $b\geq2$ and $d\geq2$. Write $[b]=\{1,\ldots,b\}$ and let
\[
 V=\bigcup_{k=0}^{d}[b]^k,
 \qquad V_{\rm int}=\bigcup_{k=0}^{d-1}[b]^k,
 \qquad V_{\rm leaf}=[b]^d.
\]
The empty word $\epsilon$ is the root. A world is a collection
\[
 \theta_v\in[b]\quad(v\in V_{\rm int}),
 \qquad z_v\in\{0,1\}\quad(v\in V_{\rm leaf}).
\]
All choices are independently possible. Starting at $v_0=\epsilon$,
define $v_{k+1}=v_k\theta_{v_k}$ and put $z=z_{v_d}$. The task is to
identify $z$ soundly and manufacture a reusable terminal device $H_z$.
On a subsequent fresh client input $Q$, that device must implement the
conversion $Q\to V_z$ while returning its installed stock. The client
input is supplied only after manufacture. The physical output species
$V_0,V_1$ are distinct from every diagnostic reporter species.

Both policies begin with the same root probe already installed and perform
the same charged root measurement, obtaining $r=\theta_\epsilon$.
This common initial evidence selects the subtree
\[
 V_r=\{rw:0\leq |w|\leq d-1\},
 \qquad S_d=|V_r|=1+b+\cdots+b^{d-1}.
\]
The comparison below starts after that common measurement. Thus the root
is not charged again as a newly manufactured probe. There are exactly $d$
further probes on the realized route and $b^{d-1}$ possible terminal leaves
conditional on $r$.

\subsection{One fixed executable, mass-conserving substrate}
There are generic assembly species $A,B,F$, retained waste $W$, and, for
every nonroot address $v$, an intermediate $D_v$ and a probe catalyst
$C_v$. The assembly instructions are exactly
\[
 A+B\longrightarrow D_v,
 \qquad D_v+F\longrightarrow C_v+W.
\tag{M$_v$}\label{ing:manufacture}
\]
There are analogous two-step instructions with $D^H_j,H_j$ in place of
$D_v,C_v$, for $j=0,1$. All these instructions are supplied from the
beginning, independently of the world and evidence. No instruction
changes one address-specific catalyst into another or returns $A,B,F$.
A manufacture budget $B_0$ means initially $B_0$ copies of each of
$A,B,F$. Every manufactured probe or terminal catalyst consumes one
batch $A+B+F$. Partial assemblies cannot increase this count.

The plant contains exactly one hidden register token $R_{v,a_v}$ at each
address, where $a_v=\theta_v$ at an internal node and $a_v=z_v$ at a leaf.
Its initial stock is
\[
 P_w=\sum_{v\in V}R_{v,a_v}.
\]
The root catalyst $C_\epsilon$ is also supplied. Register tags are not
visible to the fixed detector. For each legal $v,a$, the substrate has
exactly the following probe mechanism:
\[
 C_v+R_{v,a}+U\longrightarrow I_{v,a},\qquad
 I_{v,a}\longrightarrow J_{v,a},\qquad
 J_{v,a}\longrightarrow C_v+R_{v,a}+O_a.
\tag{P$_{v,a}$}\label{ing:probe}
\]
Here $U$ is one fresh diagnostic reporter and the distinct $O_a$ have a
fixed calibrated visible readout. An elementary instruction can be
attempted; disabledness gives a nondestructive blocked response. This
response is part of the stipulated access model and is the same for both
policy classes. The controller tries the first instruction for the legal
labels in a fixed order until one succeeds, then performs its second and
third instructions. Exactly one first instruction can succeed. This
uniform reader uses one reporter, exactly three successful firings, and
at most $|\{a:a\text{ legal at }v\}|+2$ attempts. It returns both its
catalyst and its hidden plant register exactly. The branch is learned
from the successful attempt and confirmed by its reporter; no prior
knowledge of $a$ is needed to execute the script.

The detector also observes action success, public assembly labels, and
reporter counts. It never directly observes a register tag. No other
instruction mentions $R_{v,a},I_{v,a},J_{v,a}$. In particular, without an
address-$v$ catalyst or a complex descended from it, every address-$v$
probe attempt is blocked independently of $a_v$. There is no generic
reader, free relabelling, programmable substitute catalyst, or uncharged
external instrument that reads these tags. This is an architectural
premise of the native example, not a conclusion about arbitrary physical
laboratories. The permitted controls and detector remain fixed throughout;
manufacture changes enabled native processes, not their declared access.

Reference arms use the supplied calibrated conversions $U\to O_1$ and
$U\to O_0$. A terminal catalyst has the separate service mechanism
\[
 H_j+Q\longrightarrow K_j,\qquad K_j\longrightarrow L_j,
 \qquad L_j\longrightarrow H_j+V_j.
\tag{T$_j$}\label{ing:terminal}
\]
There are no other instructions involving $Q,V_j,K_j,L_j$.

Give $A,B,F,W,U,Q,O_a,V_j,R_{v,a}$ weight one; give $D_v,C_v,D^H_j,H_j$
weight two; give $I_{v,a},J_{v,a}$ weight four; and give $K_j,L_j$ weight
three. Every displayed instruction conserves this positive integral
mass. The root probe and the entire plant are common initial resources.
Their size is exponential in $d$, and the available address-specific
reaction repertoire also grows with the tree. No uniform substrate-size
or instruction-description complexity bound is claimed. The result charges additional probe and
terminal manufacture, not the cost of constructing that common target
plant. An exponential ratio for total closed-universe material is therefore
not being claimed.

\subsection{Experiments and an explicit revision obligation}
A nominal operational hypothesis at an internal address predicts that its
probe agrees with the reference $U\to O_1$; at a leaf it predicts agreement
with $U\to O_0$. These predictions may be true or false in the supplied
world class. An actual diagnostic runs both arms on separate fresh
reporters. An unequal pair is an executable counterexample to that local
prediction. Under the independently supplied finite native model, the
active output also identifies the local table exactly: its value is
$a_v$. The next construction is selected from that revised local table.
This is a genuine change to a tested operational prediction when
$a_v$ differs from the nominal value, but it is finite parameter and
observational-presentation revision, not discovery of new explanatory
axioms.

There is also an exact connection to least observational repair. At an
installed set $S$ of probe addresses, consider all native experiments
using those probes without further manufacture. Two worlds have the same
possible observation transcripts exactly when their coordinates $a_v$
agree on $S$. Necessity follows by executing the reader at any differing
coordinate. Sufficiency follows by induction over the adaptive experiment:
all accessible instruction outcomes depend only on the common public
stock and those coordinates. Thus the observational quotient is
\[
 q_S:w\longmapsto (a_v)_{v\in S}.
\]
Installing $C_v$ for $v\notin S$ changes this to $q_{S\cup\{v\}}$.
The forgetful comparison
\[
 Q_{S\cup\{v\}}\twoheadrightarrow Q_S
\]
is a strict quotient map whenever the new coordinate is nonconstant on
the current evidence fibre. This is the least refinement supporting the
new reader: any quotient through which both the old tests and the reader
factor must distinguish every pair distinguished by their product.
The diagnostic does not factor through the old quotient on such a fibre.
These claims concern installed experimental repertoires, not the global
runtime, which already contains all assembly-and-test programs.

A single world need not refute the nominal prediction at every stage.
For $b\geq3$, however, the subfamily with all visited internal selectors
in $\{2,\ldots,b\}$ and terminal bit $1$ gives a measured disagreement
at every new diagnostic. The policy and the theorem are valid for the
full world class, not just this positive subfamily.

\subsection{The two policy classes}
An interleaved policy is one causal algorithm whose next native action is
a function of the finite observed history and its own stock. It may
alternate manufacture and probe use. All controller memory and computation
may be available to both sides; neither is charged in the following
material comparison.

A manufacture-once policy receives the same common root evidence. It then
has one fabrication phase, followed by an arbitrary adaptive experimental
phase, with no subsequent fabrication instruction. During fabrication it
may observe assembly completion and its deterministic public byproducts,
but may not interrogate the unknown plant beyond the common root evidence.
Equivalently, after that evidence all fabrication instructions precede all
new plant-interrogating instructions. Random choices may be made, and later
experiments, branching, computation, stopping, and reporter supply are
unrestricted. The terminal service occurs after both phases.

This is a stronger comparator than a fixed nonadaptive list of tests. It
is nevertheless a real restriction: it forbids responding to new plant
evidence by manufacturing another artifact. Allowing a prewritten
program to perform such later manufacture would put that program in the
interleaved class, however early the program was written. The temporal
restriction does not remove any physical action from one side's runtime.

\begin{theorem}[Exact manufacture-once gap and an executable policy]
\label{ing:gap}
For every $b\geq2$ and $d\geq2$, after the common root experiment:
\begin{enumerate}
\item One deterministic causal interleaved policy works in every world
using exactly $d+1$ assembly batches: $d$ probes and one terminal device.
Every selected diagnostic has an explicit native acquisition and execution
witness; the final manufacture is enabled in the actual post-experiment
residual stock. The terminal service uses a subsequent fresh $Q$.
\item A manufacture-once policy that succeeds in every world requires at
least $S_d+2$ assembly batches. This bound is attained: construct every
probe in $V_r$ and both terminal devices, then follow the route by adaptive
experiments. Hence the exact worst-case manufacture-once requirement is
\[
 S_d+2=\frac{b^d-1}{b-1}+2,
\]
whereas $d+1$ batches suffice with interleaving.
\item The lower bound still holds if the manufacture-once policy has
unlimited reporters, unlimited computation, arbitrarily many adaptive
experiments, and every read-only experiment physically implementable
using its installed portfolio. Randomization cannot give sure success
with a smaller pathwise manufacture budget.
\end{enumerate}
\end{theorem}
\begin{proof}
For the first claim, set $v=r$. Manufacture $C_v$ with the two instructions
in \eqref{ing:manufacture}, run its active and reference diagnostic, and
retain the observed local value. If $v$ is internal, replace $v$ by
$v a_v$ and repeat. After exactly $d$ new probes the leaf value $z$ is
known. Manufacture $H_z$ from the one reserved batch. This is one algorithm
for all worlds: its choice at every step depends only on completed visible
outcomes. It does not select a different policy after learning the world.
Each completed active probe restores its $C_v$ and $R_{v,a_v}$ exactly and
uses only its reporter; each reference arm uses one further reporter.
Therefore after $k\leq d$ new diagnoses the residual stock includes
$d+1-k$ untouched copies of each of $A,B,F$. The final manufacture is
enabled with $k=d$, with no reset, material copying, or ignored waste.
When a fresh $Q$ later arrives, \eqref{ing:terminal} implements the required
three-step stock-preserving service. On the realized route, the sequence
is evidence, manufacture, new experiment, revised local table, another
manufacture, and eventually terminal manufacture and service.

For the lower bound, condition on the root outcome $r$. Before any further
plant interrogation, a deterministic manufacture-once policy produces an
installed portfolio independent of the remaining world coordinates. At
most one address-specific probe or terminal catalyst can arise per
consumed assembly batch. Intermediates and unsuccessful fabrication do
not increase this count. We may only strengthen the policy by completing
any partially assembled probe for free and giving it ideal nondestructive
read access to every coordinate for which it has spent an addressed batch.
Let $S$ be this set of probe coordinates.

Suppose first that a leaf $v\in V_r$ is absent from $S$. Choose the ancestor
selectors so that the route reaches $v$, and choose two worlds agreeing at
all coordinates except $z_v$, which is respectively zero and one. Every
possible later observation is identical in the two worlds: no native
instruction can make the absent coordinate affect a visible outcome.
Induction on adaptive histories gives identical transcripts and outputs,
although the required bits differ. Thus every leaf is necessary.

Next suppose an internal $v\in V_r$ is absent from $S$. Choose its ancestor
selectors so the route reaches $v$. Choose two different children of $v$.
Fix all other selector coordinates so that each chosen child route reaches
a leaf, and assign those two leaves different bits. Take two worlds with
these same choices everywhere except at $v$, whose selector chooses the
first child in one world and the second in the other. All coordinates
accessible to the portfolio are identical, so all adaptive observation
transcripts are identical, while the required routed bits differ. Hence
every internal coordinate in $V_r$ is necessary too. This proves that at
least $S_d$ probe batches must have been committed.

Finally, before new observations the terminal bit is independently either
zero or one. No later instruction can create a previously unmanufactured
$H_j$: the terminal service only returns an $H_j$ that was already present,
and its client input has not yet been supplied. Thus the static portfolio
must contain both $H_0,H_1$, requiring two further batches. The construction
in the statement attains the bound.

For a randomized policy, condition on its random seed to obtain a
deterministic portfolio and strategy. The finite world family has a finite
number of possible coordinate portfolios. If every seed leaves some
necessary coordinate or terminal type absent, the preceding argument gives
failure on at least one world for that seed. A finite union argument shows
that some fixed world has positive failure probability. Thus almost-sure
success on every world also requires the stated pathwise budget. This
argument does not restrict later computational or experimental adaptivity.
\end{proof}

\paragraph{A complete resource ledger.}
The common initial root measurement uses one reporter and three successful
firings. The $d$ later two-arm diagnoses use $2d$ reporters and $4d$
successful firings. The $d+1$ assemblies use $d+1$ copies of each of
$A,B,F$, retain $d+1$ waste tokens $W$, and use $2(d+1)$ successful firings.
The subsequent terminal call consumes its external $Q$, produces $V_z$,
and uses three successful firings. Thus the complete interleaved run,
including the common initial measurement and one terminal call, has
exactly $6d+8$ successful firings. Counting blocked action attempts as
well gives at most $d(b+5)+9$ attempts: there are $d$ internal-node active
reads including the root, and one binary leaf read. All diagnostic outputs
are retained or exported through explicitly counted output ports. The
plant registers, root catalyst, route probes, and $H_z$ remain available.
The cost statements do not assert an optimal controller time or optimal
interleaved firing count among every conceivable successful policy.

\subsection{The reusable general principle and its boundary}
\begin{corollary}[Independent weighted manufacture]
Suppose the two assembly steps for a depth-$k$ probe have total charged
cost $c_k>0$, and a terminal assembly has cost $c_H>0$, with no alternative
construction, sharing, or recycling that decreases these charges. Then
the same policy has manufacture cost
\[
 \sum_{k=1}^{d}c_k+c_H,
\]
whereas the minimum sure-success manufacture-once cost is
\[
 \sum_{k=1}^{d}b^{k-1}c_k+2c_H.
\]
\end{corollary}
\begin{proof}
The adaptive policy builds one probe at each depth and one terminal device.
The lower-bound proof forces every one of the $b^{k-1}$ probes at depth
$k$ in the root-conditioned subtree and both terminal devices. Additivity
and the absence of cheaper shared constructions give the lower bound;
manufacturing those types separately attains it.
\end{proof}

The lower bound uses a general criterion rather than a feature unique to
color labels. Fix a common initial evidence fibre and independently
charged, irreversible instrument types. Suppose that, whenever an
instrument type $v$ is absent, there are two worlds in that fibre which
agree on the full transcript semantics of every experiment using the
remaining installed instruments but require different terminal answers.
Then every universally sound manufacture-once portfolio contains $v$.
If this condition holds separately for every type in a set $T$, it must
contain their union, with the corresponding additive resource cost.
A causal manufacture-and-experiment decision tree, by contrast, may pay
only for the types on the realized root-to-leaf path, plus its terminal
construction. This is proved by exactly the two-world indistinguishability
and residual-stock induction above. It applies beyond reaction systems
when actual native witnesses discharge its instrument and cost premises.

Categorically, the native arrows live between resource configurations,
with additive costs and observational histories retained. Each installed
configuration supplies its own experiment category and observational
quotient. Manufacture gives an actual configuration transition; a compiled
diagnostic supplies the new separating context; revision gives a quotient
comparison in the direction of forgetting distinctions. A policy chooses
its next native arrow over a completed evidence history. There is no need
to pretend that such a causal family is one endofunctor of a fixed
observational category. In this example the relevant quotient maps are
explicit coordinate projections, so no difficult categorical existence
claim is hidden in this description.

The finite native model is a special case of Section~\ref{sec:native-semantics}:
attempts have explicit blocked or successful outcomes, and the three-step
probe macro has a native witness with exact restoration. With a fixed
finite reporter and client bound it also has finite transition tables;
the reusable terminal-service relation treats each fresh client $Q$ as
an environmental input. The completed information state retains all
register tags as possible worlds and all material as actual stock.
The coordinate quotient describes the installed experiments, whereas
the changing evidence fibre describes what has been learned on this run.
The theorem is an additional-manufacture separation, not a claim that
the substrate discovers its own source class or semantic language.

The separation disappears as an impossibility if unrestricted additional
stock and manufacture are allowed: the finite static portfolio works.
It may also disappear or shrink if the substrate admits a cheap universal
reader, sufficiently cheap address reconfiguration, a direct terminal-bit
probe, correlated world parameters that determine missing coordinates,
or allowed manufacture during the supposedly static observation phase.
No transcript argument can restore a lower bound after its actual
indistinguishability premise has failed. No finite experiment by itself
warrants the stipulated conservation laws, exact restoration, stationarity,
calibration, or absence of side channels. Those are independent source
and instrument promises, exactly as in the previous native examples.


\section{What revision and categorical comparison do not supply}
\label{sec:revision-limits}
\subsection{A fixed completion is not a choice of the best theory}
For a fixed source algebra $X$, supplied finitary operations $\Omega$,
old quotient $q^0$, and diagnostic seeds $D$, the standard contextual
construction is
\[
 x\equiv_*y\quad\Longleftrightarrow\quad
 t(C[x])=t(C[y])
 \quad\text{for every admitted one-hole context }C
 \text{ and old or new seed }t.
\]
It is the largest $\Omega$-congruence refining the old distinctions and
diagnostic kernels. Indeed a context-closed indistinguishability relation
is a congruence by changing one argument of an operation at a time;
conversely every congruence retaining the seeds is preserved by all
contexts. Hence $X/{\equiv_*}\to Q^0$ is the coarsest sufficient
\emph{relative} refinement. The finite completion in
Section~\ref{msc:sec:completion} is its transition-system counterpart,
including enabledness and residual state. In the continuous example,
$q^*(\theta)=\theta$ is the two-point instance. Neither construction
creates a new diagnostic or proves its native execution. Those are
supplied by the preceding manufacture and access proofs.

There is generally no unique least empirical repair across changing
languages, candidate classes, loss criteria, or instrument choices. For
example take $F=k[x,y]$, an old selected proposal $J=(x,y)$, and two
currently supported exact candidates $I_1=(x)$ and $I_2=(y)$. Each retains
one old generator and retracts the other; neither contains the other.
There is no largest supported selected candidate under inclusion. Their
common skeptical commitment is $(x)\cap(y)=(xy)$, but it is neither selected
candidate and does not decide which is physically correct. Likewise two
alternative sensors can yield incomparable minimally sufficient
languages. A choice among them requires an additional objective and
resource contract.

For a fixed possible-world family, exact evidence filtering from $B$ to
$B'\subseteq B$ has
\[
 \operatorname{Th}(B)=\bigcap_{\theta\in B}\operatorname{Th}(\theta)
 \ \subseteq\
 \operatorname{Th}(B').
\]
It strengthens warranted common consequences. Retracting a selected
conjecture $J$ instead retains $J\cap\operatorname{Th}(B')$ and need not
preserve the rest of $J$. A forward interpretation from a quotient that
imposes a refuted equation cannot send that equation's two sides to
unequal revised denotations. A common-fragment span is correctly typed;
an ordinary law-preserving forward map generally is not.

\subsection{Semantic equality does not restore a native witness}
A semantic interpretation of the native execution category can identify
different histories. Even if every target arrow has a native witness,
these witnesses need not admit a composition-compatible selection. In a
free path category on a diamond
\[
 0\longrightarrow a\longrightarrow1,
 \qquad 0\longrightarrow b\longrightarrow1,
 \qquad a\ne b,
\]
the reachability quotient has one arrow $0\leq1$. An identity-on-objects
section preserving composition would have to send it both to the path
through $a$ and to the path through $b$. These are distinct native
histories. This is the established witness-forgetting obstruction from
\emph{Native Operations across Doctrines}; it remains present even when
all endpoints and their observational meanings agree.

The algorithms here do not assert such a section. They retain the
command program and actual execution evidence, and prove that its
observation-based branches compose at compatible residual states. Nor do
they derive a tensor product of devices sharing material, a distributive
law for arbitrary revision and wiring, or a contextad from a displayed
sequence of arrows. A richer categorical interpretation needs its actual
object map, native arrow map, observations, endpoint transport, and
unit/composition coherence, with higher cells when the doctrine has them.
A preservation theorem then transports a positive policy. Only reflection
of the declared baseline transports the negative comparison.

The bill-of-material result makes a more limited categorical claim that
is fully proved: the monotone additive map $C$ between thin resource
categories commutes with finite joins exactly under the stated row
condition. Its proof constructs physical setup histories in the direct
factory, rather than inferring physical implementation from the existence
of a thin categorical arrow. The finite completion separately gives an
exact quotient of a supplied deterministic transition model. Keeping
these types separate is necessary for the combined argument.



\section{Verification scope}\label{sec:verification}
The manuscript's general positive theorems are proved mathematically
above. Separate executable checks exercise bounded instances; they are
not substitutes for those proofs or physical validation. The direct
factory checks cover 86 small matrices, 11,430 join-pair cases, 1,870
stock/normal-form cases and 406 normalization steps. Seeded finite-model
checks for the completion theorem cover 100 models, 35,580 raw beliefs,
55,284 successor comparisons and 177,900 bounded-policy comparisons.
These are finite regressions, not exhaustive verification of all inputs.

A companion Lean development, \texttt{ResourceSensitiveLoop.lean}, checks
an executable deterministic partial-command semantics, composition at
actual residual states, the observation-uniform one-step obligation, and
the destructive-specimen impossibility of
Proposition~\ref{prop:destructive-diagnosis}. Its impossibility is by
induction over arbitrary finite policies, including arbitrarily many
idle steps, rather than an imposed search horizon. It also checks the
controlled repair in which the same test preserves the specimen and is
followed by the build selected by its output. The printed dependencies of its main theorems are the standard axioms
\texttt{propext} and \texttt{Quot.sound}; no admitted proof placeholders or
extra axioms are used in that development.

That formalization does not verify the bill-of-material theorem, finite
completion algorithm, continuous latch, source stationarity, interleaving
family, randomized-policy claims, or the full positive manufacture loop.
Its abstract controller is supplied and its memory is not fabricated or
charged. The continuous conclusions remain conditional analytic proofs
under the declared physical promises. The TeX source is delivered
uncompiled; no rendered-layout or PDF check is asserted.


\section{Precedents and the precise level of the result}\label{sec:prior}
The nearest precedents already contain most of the general machinery.
The purpose of this comparison is to identify the additional native
obligations, not to establish priority by absence from a short literature
list. The theorem numbers below refer to the cited versions.

\paragraph{Controller and strategy refinement.}
Alternating refinement formalizes the controller/environment quantifier
pattern; Alur et al. relate it to preservation of strategic temporal
properties \cite{Alur1998}. Reissig, Weber and Rungger's feedback
refinement relation, Definition~V.2 and Theorems~V.4--V.5 of the cited
version, gives sufficient and necessary conditions for uniform controller
reuse through the specified static quantizer; Theorem~VI.3 transfers
specification satisfaction \cite{RWR2017}. Reissig and Rungger's valuated
refinement already includes running and terminal costs, and their
Theorem~V.6 bounds the refined controller's worst-case value
\cite{RR2019}. Calbert, Girard and Jungers, Theorem~3.19 and
Corollary~3.20, characterize several simulation relations by particular
implementable control interfaces \cite{CGJ2026}. Thus neither scalar
costs, all-branch correctness, nor an architecture-sensitive converse is
new in itself. Proposition~\ref{prop:macro-lift} is explicitly refinement
reasoning. The additional concrete work is computing an exact completion
from primitive tables, and deriving the direct-factory interchange and
normalization criterion from a bill of materials.

\paragraph{Contingent experiments and construction.}
Albore, Palacios and Geffner's knowledge-level translation for contingent
planning has soundness and completeness results, Theorems~5--6, using all
initial-state tags \cite{APG2009}. Such policy trees already interleave
state-changing actions and sensing. The bounded recursion in
Theorem~\ref{msc:thm:joint} is this established style of strong planning,
with an explicitly computed reusable-service terminal condition and the
actual residual resources retained. Van den Bos and Vaandrager's
Theorem~6.1 constructs adaptive distinguishing graphs under its injective
splitting hypothesis, and their Example~7 shows why that hypothesis
matters \cite{BV2019}. Their work already treats destructive state
identification, partial availability, and visible nondeterminism.
Our destructive-specimen obstruction is a simple explanation of the
same need to reason about one real experiment, not a new principle of
adaptive testing.

\paragraph{Categorical learning and revision.}
CALF, Theorem~17 in the cited version, identifies a hypothesis with its
target under the stated reachability/observability and consistency or
closedness conditions \cite{CALF2017}. FunL*, Theorem~26, gives termination
at the minimal automaton under its factorization, power/copower and
Noetherian assumptions \cite{FunL2021}. Generalized tree CALF,
Theorems~18 and~25, uses counterexamples to force closedness/consistency
progress and proves termination under finite subobject/quotient
conditions \cite{TreeCALF2022}. These are direct precedents for
contextual refinement beyond ordinary words. CoalA's Theorems~4.25 and~5.3
refine component alphabets and reconstruct a synchronizing composite
\cite{CoalA2025}; it would be false to say all prior learning fixes every
component interface. Here the additional obligation is actual manufacture
and calibration of the newly admitted experimental interface. The finite
candidate source class and the elementary operational vocabulary are
still supplied.

Ma et al. explicitly combine typed theory edits and diagnostic experiments;
their Propositions~1--2 treat a within-class blind spot and stage-wise
rejection under the stated predictive assumptions \cite{Ma2026}.
Wang and Buehler's categorical revision framework includes preservation
maps, artifact states and gated transition composition
\cite{WB2026}. These are close comparators, not grounds for calling
experiment--theory coupling itself new. The present exact constructions
revise a source-law candidate or an operational quotient within a
supplied language. They do not prove open-ended scientific explanation or
invention of that language.

\paragraph{Histories, compositional interfaces and physical dynamics.}
Joyal, Nielsen and Winskel's Proposition~1 gives local path lifting for
open maps, and their Theorem~2 recovers strong bisimulation by spans of
open maps \cite{JNW1996}. Thus categorical history lifting is not a new
idea. Capucci and Myers' contextad results include the Kleisli/Ctx
comparison, trifunctoriality, and transfer of algebraic structure in the
cited version \cite{CM2026}; Spivak develops interaction-dependent
interface shapes \cite{Spivak2026}. Those constructions do not by
themselves supply this paper's native material histories or baseline
reflection. Feedback memory and stable dynamical state are also classical,
with Hopfield's graded-response networks an important antecedent
\cite{Hopfield1984}. The latch certificate is reused from the previous
native-growth note, not claimed as a new circuit discovery.

\paragraph{Exact dependence on the preceding project results.}
\emph{Endogenous Capability Growth} already proves the whole-runtime
boundary, robust feedback formation, its four-probe certificate, delayed
comparison, and invariant-relation reuse \cite{ProjectGrowth}.
\emph{From a Failed Operational Law to an Executable Diagnostic} already
contains effective diagnostic transport and tagged-belief bounded
policy recursion \cite{ProjectDiagnosis}. The revision--wiring note
already proves coarsest contextual repair relative to supplied source
operations and seeds \cite{ProjectRevision}. \emph{Native Operations
across Doctrines} already distinguishes semantics from native execution,
solves a finite native-table coherence problem, and proves the free-path
witness-selection obstruction \cite{ProjectNative}.

The added result is therefore a constructive synthesis with independently
checkable local conclusions: the direct-production criterion equates
join preservation, same-material advance manufacture, confluence and
universal capacity normalization; the finite native completion commutes
exactly with residual belief update and preserves/reflects policy
feasibility and optimum manufacture cost; and the physical policy uses
newly manufactured memory to obtain construction-relevant evidence and
build a separate, reusable corrector from reserved stock. The scaling
family shows that repeated interleaving can reduce additional manufacture
exponentially in depth under its declared instrument grammar. Its common
plant and signature are already exponentially large, so this is not an
overall complexity or total-material advantage. No universal new theory
of discovery or historical novelty claim follows.

\section{Conclusion}\label{sec:conclusion}
The proved loop is native and closed: an actually assembled instrument
makes a diagnostic executable; the new diagnostic shows that the old
characterization is insufficient for the enlarged questions, and its
result can refute a selected source-law conjecture;
the warranted distinction selects a further physical construction; and
the resulting device has a reusable service proof from the actual
residual state. In the continuous example the new datum is a delayed
relation, and the further artifact is a fixed labelled corrector. The
prior feedforward architecture cannot carry out the matched robust
one-trial calibration-to-construction task.

The structural and algorithmic results explain when the loop can be
compiled, checked, or replaced by manufacture in advance. Their limits
are explicit: candidate physics and command access are supplied;
finite observations certify only consequences inside warranted model
promises; controller and sensing resources cannot be hidden; and semantic
quotients do not recover discarded execution histories. The complete
policy remains inside its original native runtime. That boundary leaves
room for genuine gains in installed experiments, guaranteed construction
under uncertainty, and material efficiency, which are the gains established
here.


\begin{thebibliography}{99}
\bibitem{ProjectGrowth}
\emph{Endogenous Capability Growth through Native Interaction:
Manufacture, reusable interfaces, and subsequent experiments}.
Research manuscript, 5 November 2026.
Source: \texttt{endogenous\_capability\_growth.tex}.

\bibitem{ProjectDiagnosis}
\emph{From a Failed Operational Law to an Executable Diagnostic:
Relative separation, resource-preserving language change, and no-reset
obstructions}. Research manuscript, 5 November 2026.
Source: \texttt{executable\_diagnostics\_new.tex}.

\bibitem{ProjectRevision}
Research manuscript on operational revision and wiring, 5 November 2026.
Source: \texttt{revision\_wiring\_new.tex}.

\bibitem{ProjectNative}
\emph{Native Operations across Doctrines}.
Research manuscript, 5 November 2026.
Source: \texttt{native\_operations\_across\_doctrines.tex}.

\bibitem{Alur1998}
R. Alur, T. A. Henzinger, O. Kupferman and M. Y. Vardi.
\emph{Alternating Refinement Relations}. CONCUR 1998.
\url{https://www.cs.huji.ac.il/~ornak/publications/concur98.pdf}.

\bibitem{RWR2017}
G. Reissig, A. Weber and M. Rungger.
\emph{Feedback Refinement Relations for the Synthesis of Symbolic
Controllers}. IEEE Transactions on Automatic Control, 2017.
Cited version: arXiv:1503.03715v3.
\url{https://arxiv.org/abs/1503.03715}.

\bibitem{RR2019}
G. Reissig and M. Rungger.
\emph{Symbolic Optimal Control}. IEEE Transactions on Automatic
Control, 2019. Cited version: arXiv:1709.07333v2.
\url{https://arxiv.org/abs/1709.07333}.

\bibitem{CGJ2026}
L. Calbert, A. Girard and R. M. Jungers.
\emph{Characterizing simulation relations through control architectures
in abstraction-based control}. Cited version: arXiv:2410.06083v2, 2026.
\url{https://arxiv.org/abs/2410.06083v2}.

\bibitem{APG2009}
A. Albore, H. Palacios and H. Geffner.
\emph{A Translation-Based Approach to Contingent Planning}. IJCAI 2009.
\url{https://ijcai.org/proceedings/2009/271}.

\bibitem{BV2019}
P. van den Bos and F. Vaandrager.
\emph{State Identification for Labeled Transition Systems with Inputs
and Outputs}. Cited version: arXiv:1907.11034v2, 2019;
expanded journal version, 2021.
\url{https://arxiv.org/abs/1907.11034}.

\bibitem{CALF2017}
G. van Heerdt, M. Sammartino and A. Silva.
\emph{CALF: Categorical Automata Learning Framework}. CSL 2017.
Cited version: arXiv:1704.05676v3.
\url{https://arxiv.org/abs/1704.05676}.

\bibitem{FunL2021}
T. Colcombet, D. Petri\c{s}an and R. Stabile.
\emph{Learning Automata and Transducers: A Categorical Approach}.
CSL 2021, LIPIcs 183, article 15.
\url{https://drops.dagstuhl.de/storage/00lipics/lipics-vol183-csl2021/LIPIcs.CSL.2021.15/LIPIcs.CSL.2021.15.pdf}.

\bibitem{TreeCALF2022}
G. van Heerdt, T. Kapp\'e, J. Rot, M. Sammartino and A. Silva.
\emph{A Categorical Framework for Learning Generalised Tree Automata}.
CMCS 2022. Cited version: arXiv:2001.05786v2.
\url{https://arxiv.org/abs/2001.05786}.

\bibitem{CoalA2025}
L. Henry, M. R. Mousavi, T. Neele and M. Sammartino.
\emph{Compositional Active Learning of Synchronizing Systems Through
Automated Alphabet Refinement}. CONCUR 2025, LIPIcs 348, article 20.
\url{https://drops.dagstuhl.de/entities/document/10.4230/LIPIcs.CONCUR.2025.20}.

\bibitem{Ma2026}
S. Ma et al.
\emph{Learning to Outgrow a Theory: Experimental Discovery Beyond the
Initial Hypothesis Space}. arXiv:2610.07627, 2026.
\url{https://arxiv.org/abs/2610.07627}.

\bibitem{WB2026}
F. Y. Wang and M. J. Buehler.
\emph{Self-Revising Discovery Systems for Science: A Categorical
Framework for Agentic Artificial Intelligence}.
Cited version: arXiv:2606.01444v2, 2026.
\url{https://arxiv.org/abs/2606.01444v2}.

\bibitem{JNW1996}
A. Joyal, M. Nielsen and G. Winskel.
\emph{Bisimulation from Open Maps}. Information and Computation, 1996.
The cited theorem numbering is from BRICS report RS-94-7, 1994.
\url{https://www.brics.dk/RS/94/7/BRICS-RS-94-7.pdf}.

\bibitem{CM2026}
M. Capucci and D. J. Myers.
\emph{Contextful Arrows via Wreath Products}.
arXiv:2410.21889v2, revised 3 November 2026.
\url{https://arxiv.org/abs/2410.21889v2}.

\bibitem{Spivak2026}
D. I. Spivak.
\emph{Interactions that reshape the interfaces of the interacting parties}.
arXiv:2602.17917v1, 2026.
\url{https://arxiv.org/abs/2602.17917v1}.

\bibitem{Hopfield1984}
J. J. Hopfield.
\emph{Neurons with graded response have collective computational
properties like those of two-state neurons}.
Proceedings of the National Academy of Sciences 81(10) (1984), 3088--3092.
\url{https://doi.org/10.1073/pnas.81.10.3088}.
\end{thebibliography}



\end{document}
